% ============================================================================
% SOLUCIONES DEL MÓDULO 6
% ============================================================================
% Nota: las soluciones que son apps muestran shinyApp() en un bloque que no se
% ejecuta; la lógica se comprueba con testServer(), que sí se ejecuta.

\begin{solucion}{m6-1}
Agregamos dos entradas a la \codigo{ui} y ampliamos el cálculo del
\codigo{renderText()}. El IVA se aplica después del descuento.

\begin{rcode}
library(shiny)

ui <- fluidPage(
  titlePanel("Calculadora de precio con descuento"),
  numericInput("precio", "Precio de lista ($):", value = 500,
               min = 0, step = 10),
  numericInput("cantidad", "Cantidad de piezas:", value = 1, min = 1),
  sliderInput("descuento", "Descuento (%):", min = 0, max = 50,
              value = 10, step = 5),
  checkboxInput("iva", "Incluir IVA (16%)", value = FALSE),
  h3(textOutput("total"))
)

server <- function(input, output, session) {
  output$total <- renderText({
    req(input$precio, input$cantidad)       # sin valores, en blanco
    subtotal <- input$precio * input$cantidad * (1 - input$descuento / 100)
    total <- if (input$iva) subtotal * 1.16 else subtotal
    paste0("Total a pagar: ", scales::dollar(total))
  })
}

testServer(server, {
  session$setInputs(precio = 500, cantidad = 3, descuento = 10,
                    iva = FALSE)
  print(output$total)
  session$setInputs(iva = TRUE)
  print(output$total)
})
\end{rcode}
\begin{rsalida}
[1] "Total a pagar: $1,350"
[1] "Total a pagar: $1,566"
\end{rsalida}

Sin IVA, 3 piezas de \$500 con 10\% de descuento cuestan \$1,350; con IVA,
\$1,566, como pedía el enunciado. Para abrir la app:

\begin{rnoejecutar}
shinyApp(ui, server)
\end{rnoejecutar}
\end{solucion}

\begin{solucion}{m6-2}
Hay tres cambios en \archivo{app.R} de la app 02: (1) un argumento nuevo en
\codigo{filtrar\_ventas()}, (2) la entrada en la barra lateral y (3) pasar
la entrada al filtrar. Primero probamos la función modificada, de forma
autocontenida:

\begin{rcode}
library(tidyverse)
library(scales)
ventas <- read_csv("datasets/transacciones.csv", show_col_types = FALSE)

# (1) filtrar_ventas() con el argumento nuevo 'pagos' (resto igual)
filtrar_pagos <- function(datos, pagos = NULL) {
  if (length(pagos) > 0) {                 # vacío = todos los métodos
    datos <- filter(datos, metodo_pago %in% pagos)
  }
  datos
}
solo_efectivo <- filtrar_pagos(ventas, pagos = "Efectivo")
nrow(solo_efectivo)
dollar(sum(solo_efectivo$total_transaccion), accuracy = 1)
\end{rcode}
\begin{rsalida}
[1] 315
[1] "$303,735"
\end{rsalida}

Con solo ``Efectivo'' quedan 315 transacciones por \$303,735. Los cambios en
la app (la app completa modificada se probó arrancándola y con
\codigo{testServer()}):

\begin{rnoejecutar}
# (1) En filtrar_ventas(): nuevo argumento y nuevo if
filtrar_ventas <- function(datos, tienda = "Todas", categorias = NULL,
                           fechas = range(datos$fecha), pagos = NULL) {
  # ... filtros de fecha, tienda y categoría como antes ...
  if (length(pagos) > 0) {
    resultado <- filter(resultado, metodo_pago %in% pagos)
  }
  resultado
}

# (2) En sidebarPanel(), después del selectizeInput de categorías:
checkboxGroupInput("pagos", "Método de pago:",
                   choices = sort(unique(ventas$metodo_pago)),
                   selected = sort(unique(ventas$metodo_pago))),

# (3) En el server:
datos_filtrados <- reactive({
  req(input$fechas)
  filtrar_ventas(ventas, input$tienda, input$categorias, input$fechas,
                 input$pagos)
})
\end{rnoejecutar}

Si el usuario desmarca todas las casillas, \codigo{input\$pagos} vale
\codigo{NULL} y, por nuestra convención, se muestran todos los métodos.
\end{solucion}

\begin{solucion}{m6-3}
Primero la lógica. El margen bruto de cada transacción es lo cobrado menos
el costo de las piezas vendidas (\codigo{costo * cantidad}):

\begin{rcode}
library(tidyverse)
library(scales)
productos <- read_csv("datasets/productos.csv", show_col_types = FALSE)
datos <- read_csv("datasets/transacciones.csv", show_col_types = FALSE) %>%
  left_join(select(productos, producto_id, categoria, costo),
            by = "producto_id")

calcular_margen <- function(datos) {
  datos %>%
    summarise(ventas = sum(total_transaccion),
              margen = sum(total_transaccion - costo * cantidad),
              pct_margen = margen / ventas)
}

calcular_margen(datos)
datos %>%
  group_by(categoria) %>%
  calcular_margen() %>%
  arrange(pct_margen) %>%
  head(3)
\end{rcode}
\begin{rsalida}
# A tibble: 1 x 3
   ventas  margen pct_margen
    <dbl>   <dbl>      <dbl>
1 959993. 444637.      0.463
# A tibble: 3 x 4
  categoria  ventas margen pct_margen
  <chr>       <dbl>  <dbl>      <dbl>
1 Ropa       30048.  5835.      0.194
2 Hogar      63407. 13079.      0.206
3 Belleza   182143. 56054.      0.308
\end{rsalida}

En total la empresa obtiene un margen bruto de 46.3\% (\$444,637). Las
categorías con menor margen son Ropa (19.4\%), Hogar (20.6\%) y Belleza
(30.8\%): Belleza es la tercera más baja porque varios de sus productos tienen
un costo superior a su precio de catálogo (el hallazgo de calidad de datos
que ya conoces), aunque sus precios de venta reales lo compensan en parte. Justo para eso sirve el color condicional. En la
app, agregamos \codigo{calcular\_margen()} a
\archivo{R/funciones\_dashboard.R} y estas piezas a \archivo{app.R}:

\begin{rnoejecutar}
# ui: nueva fila dentro de tabItem(tabName = "resumen")
fluidRow(
  valueBoxOutput("kpi_margen", width = 6),
  valueBoxOutput("kpi_pct_margen", width = 6)
)

# server
margen <- reactive(calcular_margen(datos_filtrados()))

output$kpi_margen <- renderValueBox({
  m <- margen()$margen
  valueBox(formato_pesos(m), "Margen bruto",
           icon = icon("sack-dollar"),
           color = if (isTRUE(m >= 0)) "green" else "red")
})
output$kpi_pct_margen <- renderValueBox({
  valueBox(formato_pct(margen()$pct_margen), "% de margen",
           icon = icon("percent"), color = "blue")
})
\end{rnoejecutar}

\codigo{isTRUE()} evita un error cuando el filtro no deja ventas (el margen
sería 0 y el porcentaje \codigo{NaN}). Con la categoría Belleza y la tienda
correspondiente puedes llegar a ver la caja en rojo.
\end{solucion}

\begin{solucion}{m6-4}
El truco es que \emph{todas} las salidas dependan de un
\codigo{eventReactive()} ligado al botón, y no directamente de los filtros.

\begin{rcode}
library(shiny)
library(tidyverse)
library(scales)
tiendas <- read_csv("datasets/tiendas.csv", show_col_types = FALSE)
ventas <- read_csv("datasets/transacciones.csv", show_col_types = FALSE) %>%
  left_join(select(tiendas, tienda_id, nombre_tienda), by = "tienda_id")

ui <- fluidPage(
  titlePanel("Ventas con botón Aplicar"),
  sidebarLayout(
    sidebarPanel(
      selectizeInput("tiendas", "Tiendas:",
                     choices = sort(unique(ventas$nombre_tienda)),
                     multiple = TRUE, options = list(placeholder = "Todas")),
      dateRangeInput("fechas", "Periodo:", start = "2023-01-01",
                     end = "2023-12-31", language = "es",
                     separator = " a "),
      actionButton("aplicar", "Aplicar", icon = icon("filter"))
    ),
    mainPanel(h3(textOutput("ventas")), h4(textOutput("transacciones")))
  )
)

server <- function(input, output, session) {
  # ignoreNULL = FALSE: calcula también al abrir la app (con el valor 0
  # del botón), para no empezar con la pantalla vacía
  datos <- eventReactive(input$aplicar, {
    d <- filter(ventas, fecha >= input$fechas[1], fecha <= input$fechas[2])
    if (length(input$tiendas) > 0) {
      d <- filter(d, nombre_tienda %in% input$tiendas)
    }
    d
  }, ignoreNULL = FALSE)

  output$ventas <- renderText(
    paste("Ventas:", dollar(sum(datos()$total_transaccion), accuracy = 1)))
  output$transacciones <- renderText(
    paste("Transacciones:", nrow(datos())))
}

testServer(server, {
  session$setInputs(tiendas = NULL, aplicar = 0,
                    fechas = as.Date(c("2023-01-01", "2023-12-31")))
  print(output$ventas)
  session$setInputs(tiendas = "Sucursal Norte")      # sin pulsar
  print(output$ventas)
  session$setInputs(aplicar = 1)                     # pulsa Aplicar
  print(output$ventas)
  print(output$transacciones)
})
\end{rcode}
\begin{rsalida}
[1] "Ventas: $959,993"
[1] "Ventas: $959,993"
[1] "Ventas: $89,024"
[1] "Transacciones: 99"
\end{rsalida}

La prueba lo demuestra: al elegir ``Sucursal Norte'' sin pulsar, las ventas
siguen en \$959,993; al pulsar ``Aplicar'' cambian a las de la Sucursal
Norte (\$89,024 en 99 transacciones).

\begin{rnoejecutar}
shinyApp(ui, server)
\end{rnoejecutar}
\end{solucion}

\begin{solucion}{m6-5}
Primero la función de la gráfica. Con \codigo{wday(fecha, label = TRUE,
week\_start = 1)} obtenemos el día como factor ordenado de lunes a domingo
(en español, porque trabajamos con locale es\_MX):

\begin{rcode}
library(tidyverse)
library(lubridate)
library(scales)
datos <- read_csv("datasets/transacciones.csv", show_col_types = FALSE)

ventas_por_dia <- function(datos) {
  datos %>%
    mutate(dia = wday(fecha, label = TRUE, abbr = FALSE,
                      week_start = 1)) %>%
    group_by(dia) %>%
    summarise(ventas = sum(total_transaccion), .groups = "drop") %>%
    mutate(es_mejor = ventas == max(ventas))
}

grafica_dias <- function(datos) {
  ggplot(ventas_por_dia(datos), aes(x = dia, y = ventas, fill = es_mejor)) +
    geom_col(width = 0.7) +
    scale_fill_manual(values = c(`TRUE` = "#eb6834", `FALSE` = "#b5b3ad"),
                      guide = "none") +
    scale_y_continuous(labels = label_dollar(),
                       expand = expansion(mult = c(0, 0.05))) +
    labs(x = NULL, y = "Ventas") +
    theme_minimal(base_size = 12) +
    theme(panel.grid.minor = element_blank())
}

ventas_por_dia(datos)
p_dias <- grafica_dias(datos)
\end{rcode}
\begin{rsalida}
# A tibble: 7 x 3
  dia        ventas es_mejor
  <ord>       <dbl> <lgl>   
1 lunes     128878. FALSE   
2 martes    119364. FALSE   
3 miércoles 141329. FALSE   
4 jueves    145280. FALSE   
5 viernes   161266. TRUE    
6 sábado    115172. FALSE   
7 domingo   148703. FALSE   
\end{rsalida}

En todo 2023 el mejor día es el viernes (\$161,266) y el más flojo el
sábado (\$115,172). Como la función recibe los datos, en la app basta con pasarle
\codigo{datos\_validos()} para que respete los filtros globales. Cambios en
la app 03:

\begin{rnoejecutar}
# En sidebarMenu(), después de "Clientes":
menuItem("Días", tabName = "dias", icon = icon("calendar-week")),

# En tabItems(), una pestaña nueva:
tabItem(tabName = "dias",
  fluidRow(
    box(title = "Ventas por día de la semana", width = 12,
        status = "primary", plotOutput("graf_dias", height = 380))
  )
),

# En el server:
output$graf_dias <- renderPlot({
  grafica_dias(datos_validos())
}, res = 96)
\end{rnoejecutar}

La función \codigo{grafica\_dias()} (con \codigo{ventas\_por\_dia()}) va en
\archivo{R/funciones\_dashboard.R}, junto a las demás gráficas.
\end{solucion}

\begin{solucion}{m6-6}
Para que la descarga contenga ``exactamente la tabla que se ve'', guardamos
la tabla en un reactivo que usan \emph{tanto} la tabla como los dos botones.
Así no hay forma de que diverjan. Probemos primero, autocontenido, lo que
escribirá cada botón:

\begin{rcode}
library(tidyverse)
library(scales)
library(writexl)
productos <- read_csv("datasets/productos.csv", show_col_types = FALSE)
ventas <- read_csv("datasets/transacciones.csv", show_col_types = FALSE) %>%
  left_join(select(productos, producto_id, nombre_producto, categoria),
            by = "producto_id")

top_productos <- function(datos, n = 10) {
  datos %>%
    group_by(Producto = nombre_producto, `Categoría` = categoria) %>%
    summarise(Unidades = sum(cantidad),
              Ventas   = sum(total_transaccion), .groups = "drop") %>%
    slice_max(Ventas, n = n, with_ties = FALSE) %>%
    mutate(Ventas = dollar(Ventas))
}

dir.create("resultados", showWarnings = FALSE)
nombre <- paste0("top_productos_", Sys.Date())
top5 <- top_productos(ventas, n = 5)
write_excel_csv(top5, file.path("resultados", paste0(nombre, ".csv")))
write_xlsx(top5, file.path("resultados", paste0(nombre, ".xlsx")))
read_csv(file.path("resultados", paste0(nombre, ".csv")),
         show_col_types = FALSE)
\end{rcode}
\begin{rsalida}
# A tibble: 5 x 4
  Producto    Categoría  Unidades Ventas    
  <chr>       <chr>         <dbl> <chr>     
1 Producto 06 Juguetes         82 $34,028.04
2 Producto 19 Deportes         85 $28,367.95
3 Producto 31 Alimentos        74 $25,429.70
4 Producto 47 Automotriz       82 $24,527.21
5 Producto 46 Automotriz       74 $24,470.29
\end{rsalida}

Los cambios en la app 02:

\begin{rnoejecutar}
# ui: la pestaña "Top productos" ahora tiene los botones arriba
tabPanel("Top productos",
         br(),
         downloadButton("bajar_csv", "Descargar CSV"),
         downloadButton("bajar_excel", "Descargar Excel"),
         br(), br(),
         tableOutput("tabla_top"))

# server: un reactivo con la tabla visible
tabla_visible <- reactive(top_productos(datos_filtrados(), n = input$n_top))

output$tabla_top <- renderTable(tabla_visible(), align = "llrr",
                                digits = 0)

output$bajar_csv <- downloadHandler(
  filename = function() paste0("top_productos_", Sys.Date(), ".csv"),
  content  = function(file) write_excel_csv(tabla_visible(), file)
)
output$bajar_excel <- downloadHandler(
  filename = function() paste0("top_productos_", Sys.Date(), ".xlsx"),
  content  = function(file) writexl::write_xlsx(tabla_visible(), file)
)
\end{rnoejecutar}

Para probarlo con \codigo{testServer()} sobre la app modificada, basta con
leer la salida del botón, que devuelve la ruta del archivo generado:

\begin{rnoejecutar}
testServer(shinyAppDir("mis-apps/02-mi-explorador"), {
  session$setInputs(tienda = "Todas", categorias = NULL, n_top = 5,
                    fechas = as.Date(c("2023-01-01", "2023-12-31")),
                    metrica = "ventas")
  read_csv(output$bajar_csv, show_col_types = FALSE)   # 5 filas
})
\end{rnoejecutar}
\end{solucion}

\begin{solucion}{m6-7}
El módulo recibe los datos como \emph{reactivo} (sin paréntesis al
pasarlo) y el nombre de la columna como texto. Dentro, cada id pasa por
\codigo{ns()}, así que las dos copias no chocan aunque ambas tengan un
deslizador llamado \codigo{"n"}.

\begin{rcode}
library(shiny)
library(shinydashboard)
library(tidyverse)
library(scales)
tiendas <- read_csv("datasets/tiendas.csv", show_col_types = FALSE)
productos <- read_csv("datasets/productos.csv", show_col_types = FALSE)
ventas <- read_csv("datasets/transacciones.csv", show_col_types = FALSE) %>%
  left_join(select(tiendas, tienda_id, ciudad), by = "tienda_id") %>%
  left_join(select(productos, producto_id, nombre_producto),
            by = "producto_id")

# ---- Módulo ------------------------------------------------------------
ranking_ui <- function(id, titulo) {
  ns <- NS(id)
  box(title = titulo, width = 6, status = "primary",
      sliderInput(ns("n"), "Mostrar top:", min = 3, max = 10, value = 5),
      plotOutput(ns("grafica"), height = 300))
}

ranking_server <- function(id, datos, columna) {
  moduleServer(id, function(input, output, session) {
    resumen <- reactive({
      datos() %>%
        group_by(nombre = .data[[columna]]) %>%
        summarise(ventas = sum(total_transaccion), .groups = "drop") %>%
        slice_max(ventas, n = input$n, with_ties = FALSE)
    })
    output$grafica <- renderPlot({
      ggplot(resumen(), aes(x = ventas, y = reorder(nombre, ventas))) +
        geom_col(fill = "#2a78d6", width = 0.7) +
        scale_x_continuous(labels = label_dollar()) +
        labs(x = "Ventas", y = NULL) +
        theme_minimal(base_size = 12)
    }, res = 96)
    resumen                    # el módulo devuelve su tabla (reactivo)
  })
}

# ---- App que usa el módulo dos veces -----------------------------------
ui <- dashboardPage(
  dashboardHeader(title = "Rankings"),
  dashboardSidebar(disable = TRUE),
  dashboardBody(fluidRow(
    ranking_ui("productos", "Top productos"),
    ranking_ui("ciudades", "Top ciudades")
  ))
)
server <- function(input, output, session) {
  datos <- reactive(ventas)
  ranking_server("productos", datos, "nombre_producto")
  ranking_server("ciudades", datos, "ciudad")
}

# Cada copia del módulo se prueba por separado con su propio 'n'
testServer(ranking_server,
           args = list(datos = reactive(ventas), columna = "ciudad"), {
  session$setInputs(n = 3)
  print(resumen())
})
testServer(ranking_server,
           args = list(datos = reactive(ventas),
                       columna = "nombre_producto"), {
  session$setInputs(n = 4)
  print(nrow(resumen()))
})
\end{rcode}
\begin{rsalida}
# A tibble: 3 x 2
  nombre       ventas
  <chr>         <dbl>
1 CDMX        383588.
2 Guadalajara 169260.
3 Tijuana     114377.
[1] 4
\end{rsalida}

La copia de ciudades con \codigo{n = 3} devuelve las tres ciudades con más
ventas (la CDMX, con cuatro tiendas, domina con \$383,588) y la de productos
con \codigo{n = 4} devuelve 4 filas: cada módulo responde a su propio
deslizador. En la \codigo{ui} completa los ids reales son
\codigo{productos-n} y \codigo{ciudades-n}.

\begin{rnoejecutar}
shinyApp(ui, server)
\end{rnoejecutar}
\end{solucion}

\begin{solucion}{m6-8}
Seguimos la receta del módulo: datos fuera del \codigo{server}, funciones
de lógica probadas por separado y un servidor corto.

\begin{rcode}
library(shiny)
library(shinydashboard)
library(tidyverse)
library(scales)
library(DT)
library(writexl)

empleados <- read_csv("datasets/empleados.csv", show_col_types = FALSE)

# ---- Lógica --------------------------------------------------------------
filtrar_empleados <- function(datos, deptos = NULL, sucursales = NULL) {
  if (length(deptos) > 0) datos <- filter(datos, departamento %in% deptos)
  if (length(sucursales) > 0) {
    datos <- filter(datos, sucursal %in% sucursales)
  }
  datos
}

kpis_rrhh <- function(datos) {
  datos %>%
    summarise(plantilla = n(),
              salario_prom = mean(salario_mensual),
              antiguedad_prom = mean(antiguedad_anos),
              pct_posgrado = mean(nivel_educacion %in%
                                    c("Maestría", "Doctorado")))
}

grafica_salario_puesto <- function(datos) {
  datos %>%
    group_by(puesto) %>%
    summarise(salario = mean(salario_mensual), .groups = "drop") %>%
    ggplot(aes(x = salario, y = reorder(puesto, salario))) +
    geom_col(fill = "#2a78d6", width = 0.7) +
    scale_x_continuous(labels = label_dollar(),
                       expand = expansion(mult = c(0, 0.05))) +
    labs(x = "Salario mensual promedio", y = NULL) +
    theme_minimal(base_size = 12) +
    theme(panel.grid.minor = element_blank())
}

tabla_empleados <- function(datos) {
  datos %>%
    select(Nombre = nombre, Departamento = departamento, Puesto = puesto,
           Sucursal = sucursal, Salario = salario_mensual,
           `Antigüedad` = antiguedad_anos) %>%
    arrange(desc(Salario))
}

kpis_rrhh(empleados)
kpis_rrhh(filtrar_empleados(empleados, deptos = "Ventas"))
\end{rcode}
\begin{rsalida}
# A tibble: 1 x 4
  plantilla salario_prom antiguedad_prom pct_posgrado
      <int>        <dbl>           <dbl>        <dbl>
1       100       27481.            3.96         0.35
# A tibble: 1 x 4
  plantilla salario_prom antiguedad_prom pct_posgrado
      <int>        <dbl>           <dbl>        <dbl>
1        16       29032.            4.14        0.312
\end{rsalida}

La empresa tiene 100 empleados con salario mensual promedio de \$27,481 y
casi 4 años de antigüedad promedio; el 35\% tiene maestría o doctorado. El
departamento de Ventas tiene 16 personas, con salario promedio de \$29,032. Ahora la interfaz y el servidor:

\begin{rcode}
ui <- dashboardPage(
  dashboardHeader(title = "Recursos Humanos"),
  dashboardSidebar(
    selectizeInput("deptos", "Departamentos:",
                   choices = sort(unique(empleados$departamento)),
                   multiple = TRUE, options = list(placeholder = "Todos")),
    selectizeInput("sucursales", "Sucursales:",
                   choices = sort(unique(empleados$sucursal)),
                   multiple = TRUE, options = list(placeholder = "Todas"))
  ),
  dashboardBody(
    fluidRow(
      valueBoxOutput("plantilla", width = 3),
      valueBoxOutput("salario", width = 3),
      valueBoxOutput("antiguedad", width = 3),
      valueBoxOutput("posgrado", width = 3)
    ),
    fluidRow(
      box(title = "Salario promedio por puesto", width = 5,
          status = "primary", plotOutput("graf_puesto", height = 320)),
      box(title = "Empleados", width = 7, status = "primary",
          downloadButton("bajar", "Descargar Excel"), br(), br(),
          DTOutput("tabla"))
    )
  )
)

server <- function(input, output, session) {
  datos <- reactive(filtrar_empleados(empleados, input$deptos,
                                      input$sucursales))
  kpis <- reactive(kpis_rrhh(datos()))

  output$plantilla <- renderValueBox(
    valueBox(kpis()$plantilla, "Plantilla", icon = icon("users")))
  output$salario <- renderValueBox(
    valueBox(dollar(kpis()$salario_prom, accuracy = 1),
             "Salario mensual promedio", icon = icon("money-bill")))
  output$antiguedad <- renderValueBox(
    valueBox(paste(round(kpis()$antiguedad_prom, 1), "años"),
             "Antigüedad promedio", icon = icon("clock")))
  output$posgrado <- renderValueBox(
    valueBox(percent(kpis()$pct_posgrado, accuracy = 1),
             "Con maestría o doctorado", icon = icon("graduation-cap")))

  output$graf_puesto <- renderPlot({
    validate(need(nrow(datos()) > 0, "No hay empleados con esos filtros."))
    grafica_salario_puesto(datos())
  }, res = 96)

  output$tabla <- renderDT(
    datatable(tabla_empleados(datos()), rownames = FALSE,
              options = list(pageLength = 8)) %>%
      formatCurrency("Salario", currency = "$", digits = 0))

  output$bajar <- downloadHandler(
    filename = function() paste0("empleados_", Sys.Date(), ".xlsx"),
    content  = function(file) write_xlsx(tabla_empleados(datos()), file)
  )
}

testServer(server, {
  session$setInputs(deptos = c("IT", "Finanzas"), sucursales = NULL)
  cat("Plantilla IT + Finanzas:", kpis()$plantilla, "\n")
  cat("Filas en el Excel:",
      nrow(readxl::read_excel(output$bajar)), "\n")
})
\end{rcode}
\begin{rsalida}
Plantilla IT + Finanzas: 18 
Filas en el Excel: 18 
\end{rsalida}

La prueba confirma que los filtros múltiples funcionan y que el Excel
descargado contiene exactamente a los empleados filtrados. Para abrir el
dashboard:

\begin{rnoejecutar}
shinyApp(ui, server)
\end{rnoejecutar}
\end{solucion}
